\documentclass{article}
\usepackage{amsmath,amssymb,amsthm}
\usepackage{algorithm,algorithmic}
\newtheorem{theorem}{Theorem}
\newtheorem{lemma}{Lemma}
\newtheorem{assumption}{Assumption}
\newtheorem{definition}{Definition}
\newcommand{\R}{\mathbb{R}}
\newcommand{\E}{\mathbb{E}}

\begin{document}

\section*{Algorithm Name}
\textbf{Accelerated Gradient Method (AGM) for $\mu$–strongly convex, $L$–smooth functions}\\
The scheme is the classical Nesterov acceleration with a \emph{fixed} momentum coefficient
\[
\beta \;=\;\frac{\sqrt{L}-\sqrt{\mu}}{\sqrt{L}+\sqrt{\mu}}\in[0,1),
\qquad 
\alpha \;=\;\frac{1}{L}>0 .
\]

\section*{Mathematical Formulation}
Let $f:\R^{d}\rightarrow\R$ be differentiable.  
Initialize $x_{0}=y_{0}\in\R^{d}$ and set $x_{-1}=x_{0}$.  
For $k=0,1,2,\dots$ perform
\begin{align}
y_{k} &= x_{k}+ \beta\,(x_{k}-x_{k-1}), \label{eq:y}\\[4pt]
x_{k+1} &= y_{k} - \alpha \,\nabla f(y_{k}). \label{eq:x}
\end{align}
The pair $(\alpha,\beta)$ is chosen as above and stays constant throughout the iterations.

\section*{Pseudocode}
\begin{algorithm}[H]
\caption{Accelerated Gradient Method (AGM)}
\begin{algorithmic}[1]
\REQUIRE $f$, $L$, $\mu$, initial point $x_{0}\in\R^{d}$, number of iterations $T$  
\STATE Set $\alpha \gets 1/L$, $\beta \gets (\sqrt{L}-\sqrt{\mu})/(\sqrt{L}+\sqrt{\mu})$  
\STATE $x_{-1}\gets x_{0}$, $y_{0}\gets x_{0}$  
\FOR{$k=0$ \TO $T-1$}
    \STATE $y_{k}\gets x_{k}+ \beta\,(x_{k}-x_{k-1})$ \hfill\textcolor{gray}{\# Eq.~\eqref{eq:y}}
    \STATE $g_{k}\gets \nabla f(y_{k})$
    \STATE $x_{k+1}\gets y_{k} - \alpha\, g_{k}$ \hfill\textcolor{gray}{\# Eq.~\eqref{eq:x}}
\ENDFOR
\RETURN $x_{T}$
\end{algorithmic}
\end{algorithm}

\section*{Dry Run Example}
Consider the univariate quadratic
\[
f(w)= (w-3)^{2},\qquad \nabla f(w)=2(w-3).
\]
It is $L$–smooth with $L=2$ and $\mu$–strongly convex with $\mu=2$.  
Hence $\beta=0$ and $\alpha=1/L=0.5$.

\begin{center}
\begin{tabular}{c|c|c|c|c}
$k$ & $x_{k}$ & $y_{k}=x_{k}+\beta(x_{k}-x_{k-1})$ & $g_{k}=\nabla f(y_{k})$ & $x_{k+1}=y_{k}-\alpha g_{k}$ \\ \hline
0 & $0$ & $0$ & $2(0-3)=-6$ & $0-0.5(-6)=3$ \\
1 & $3$ & $3$ & $2(3-3)=0$ & $3-0.5(0)=3$ \\
2 & $3$ & $3$ & $0$ & $3$ \\
\end{tabular}
\end{center}
The objective values are $f(x_{0})=9$, $f(x_{1})=0$, $f(x_{2})=0$, confirming that the method reaches the minimizer $w^{*}=3$ in one step for this simple quadratic.

\newpage
\section*{Assumptions}
\begin{assumption}[Smoothness]\label{ass:smooth}
$f$ is $L$‑smooth, i.e.
\[
\|\nabla f(u)-\nabla f(v)\|\le L\|u-v\|,\qquad\forall u,v\in\R^{d}.
\tag{A1}
\]
Equivalently,
\[
f(v)\le f(u)+\langle\nabla f(u),v-u\rangle+\frac{L}{2}\|v-u\|^{2}.
\tag{A1'}
\]
\end{assumption}

\begin{assumption}[Strong Convexity]\label{ass:strong}
$f$ is $\mu$‑strongly convex, i.e.
\[
f(v)\ge f(u)+\langle\nabla f(u),v-u\rangle+\frac{\mu}{2}\|v-u\|^{2},
\qquad\forall u,v\in\R^{d}.
\tag{A2}
\]
\end{assumption}

\begin{assumption}[Parameter Choice]\label{ass:param}
The step size and momentum are fixed as
\[
\alpha=\frac{1}{L},\qquad
\beta=\frac{\sqrt{L}-\sqrt{\mu}}{\sqrt{L}+\sqrt{\mu}}.
\tag{A3}
\]
\end{assumption}

All three assumptions hold throughout the analysis.

\section*{Main Convergence Theorem}
\begin{theorem}[Linear convergence of AGM]\label{thm:linear}
Let $f$ satisfy Assumptions~\ref{ass:smooth}--\ref{ass:strong} and let the iterates $\{x_{k}\}$ be generated by \eqref{eq:y}--\eqref{eq:x} with parameters \eqref{eq:y}--\eqref{eq:x}. Then for every $k\ge0$
\[
f(x_{k})-f(x^{\star})\;\le\;\bigl(1-\sqrt{\mu/L}\,\bigr)^{k}\,\bigl(f(x_{0})-f(x^{\star})\bigr),
\tag{T1}
\]
where $x^{\star}= \arg\min f$. Consequently,
\[
\|x_{k}-x^{\star}\|\;\le\;\Bigl(1-\sqrt{\mu/L}\Bigr)^{k/2}\,
\sqrt{\frac{2}{\mu}\bigl(f(x_{0})-f(x^{\star})\bigr)}.
\tag{T2}
\]
\end{theorem}

\section*{Theoretical Properties}
\begin{itemize}
\item \textbf{Stability.} The spectral radius of the linear operator governing the error dynamics equals $1-\sqrt{\mu/L}<1$, guaranteeing asymptotic stability.
\item \textbf{Hyper‑parameter sensitivity.} The optimal $\beta$ in \eqref{ass:param} is the unique value that minimizes the contraction factor $1-\sqrt{\mu/L}$. Any deviation $\tilde\beta\neq\beta$ yields a larger factor $1-\theta$ with $\theta<\sqrt{\mu/L}$.
\item \textbf{Comparison to baselines.}
    \begin{itemize}
        \item Gradient descent with step $1/L$ enjoys $f(x_{k})-f(x^{\star})\le (1-\mu/L)^{k}$, a slower rate because $1-\mu/L > 1-\sqrt{\mu/L}$ for $0<\mu<L$.
        \item Heavy‑ball method attains the same asymptotic factor only under a restrictive \emph{spectral} condition; AGM works for any $L$‑smooth $\mu$‑strongly convex $f$.
    \end{itemize}
\end{itemize}

\section*{Discussion}
The theorem shows that AGM achieves the optimal linear rate for first‑order methods on the class of $\mu$‑strongly convex, $L$‑smooth functions. The proof relies only on elementary inequalities derived from smoothness and strong convexity, thus illustrating that the accelerated scheme does not need any stochastic assumptions.

\newpage
\section*{Preliminaries (Definitions and Lemmas)}
\begin{definition}[Error vectors]
Denote the optimal point $x^{\star}$ and define
\[
e_{k}\;:=\;x_{k}-x^{\star},\qquad
s_{k}\;:=\;y_{k}-x^{\star}.
\]
\end{definition}

\begin{lemma}[Three‑point identity]\label{lem:three}
For any $u,v,w\in\R^{d}$ and any $\lambda>0$,
\[
\|u-w\|^{2}= \|v-w\|^{2}+2\langle u-v,\,v-w\rangle+\|u-v\|^{2}.
\tag{L1}
\]
\end{lemma}
\begin{proof}
Expand $\|u-w\|^{2}=\langle u-w,u-w\rangle$ and regroup terms. \hfill $\square$
\end{proof}

\begin{lemma}[Descent Lemma]\label{lem:descent}
If $f$ is $L$‑smooth then for any $z$,
\[
f(z-\alpha\nabla f(z))\;\le\;f(z)-\frac{\alpha}{2}\,\|\nabla f(z)\|^{2},
\qquad\text{for }\alpha\le\frac{1}{L}.
\tag{L2}
\]
\end{lemma}
\begin{proof}
From smoothness (A1') with $u=z$, $v=z-\alpha\nabla f(z)$,
\[
f(z-\alpha\nabla f(z))\le f(z)-\alpha\|\nabla f(z)\|^{2}
+\frac{L\alpha^{2}}{2}\|\nabla f(z)\|^{2}
= f(z)-\frac{\alpha}{2}\|\nabla f(z)\|^{2},
\]
since $\alpha\le1/L$. \hfill $\square$
\end{proof}

\section*{Main Proof (Step‑by‑Step Derivation)}
We prove Theorem~\ref{thm:linear}.  
All derivations are deterministic, thus expectations are omitted.

\subsection*{Step 1 – Relating $s_{k}$ and $e_{k}$}
From the definition of $y_{k}$ (Eq.~\eqref{eq:y}) and the error vectors,
\[
s_{k}=y_{k}-x^{\star}=x_{k}+ \beta(x_{k}-x_{k-1})-x^{\star}
= e_{k}+ \beta(e_{k}-e_{k-1}). \tag{S1}
\]

\subsection*{Step 2 – Update of $e_{k+1}$}
Using Eq.~\eqref{eq:x},
\[
e_{k+1}=x_{k+1}-x^{\star}=y_{k}-\alpha\nabla f(y_{k})-x^{\star}
= s_{k}-\alpha\nabla f(s_{k}+x^{\star}). \tag{S2}
\]

\subsection*{Step 3 – Quadratic Lyapunov function}
Define the potential
\[
\Phi_{k}\;:=\;\|e_{k}\|^{2}+c\,\|e_{k}-e_{k-1}\|^{2},
\qquad c>0\ \text{to be chosen}. \tag{S3}
\]
Our goal is to find $c$ such that $\Phi_{k+1}\leq\rho\,\Phi_{k}$ with $\rho=1-\sqrt{\mu/L}$.

\subsection*{Step 4 – Expand $\|e_{k+1}\|^{2}$}
Using Lemma~\ref{lem:three} with $u=e_{k+1}$, $v=s_{k}$, $w=0$,
\[
\|e_{k+1}\|^{2}= \|s_{k}\|^{2}
-2\alpha\langle s_{k},\nabla f(y_{k})\rangle
+\alpha^{2}\|\nabla f(y_{k})\|^{2}. \tag{S4}
\]

\subsection*{Step 5 – Apply strong convexity at $y_{k}$}
From (A2) with $u=y_{k}$, $v=x^{\star}$,
\[
\langle\nabla f(y_{k}),\,s_{k}\rangle
\;\ge\; f(y_{k})-f(x^{\star})+\frac{\mu}{2}\|s_{k}\|^{2}. \tag{S5}
\]

\subsection*{Step 6 – Apply smoothness to bound $\|\nabla f(y_{k})\|^{2}$}
From (A1') with $u=y_{k}$, $v=x^{\star}$,
\[
f(y_{k})-f(x^{\star})\le 
\langle\nabla f(y_{k}),s_{k}\rangle-\frac{1}{2L}\|\nabla f(y_{k})\|^{2},
\]
which rearranges to
\[
\|\nabla f(y_{k})\|^{2}\le 2L\bigl(\langle\nabla f(y_{k}),s_{k}\rangle-
f(y_{k})+f(x^{\star})\bigr). \tag{S6}
\]

\subsection*{Step 7 – Substitute (S5) and (S6) into (S4)}
Insert the lower bound (S5) for the inner product and the upper bound (S6) for the norm:
\[
\begin{aligned}
\|e_{k+1}\|^{2}
&\le \|s_{k}\|^{2}
-2\alpha\Bigl(f(y_{k})-f(x^{\star})+\frac{\mu}{2}\|s_{k}\|^{2}\Bigr)\\
&\qquad +\alpha^{2}\,2L\Bigl(\langle\nabla f(y_{k}),s_{k}\rangle-
f(y_{k})+f(x^{\star})\Bigr). 
\end{aligned}
\tag{S7}
\]

\subsection*{Step 8 – Collect terms}
Group the coefficients of $\|s_{k}\|^{2}$ and $f(y_{k})-f(x^{\star})$:
\[
\|e_{k+1}\|^{2}
\le\bigl(1-\alpha\mu\bigr)\|s_{k}\|^{2}
-2\alpha\bigl(1-\alpha L\bigr)\bigl(f(y_{k})-f(x^{\star})\bigr)
+2\alpha^{2}L\langle\nabla f(y_{k}),s_{k}\rangle . \tag{S8}
\]

\subsection*{Step 9 – Eliminate the remaining inner product}
From Lemma~\ref{lem:three} with $u=s_{k}$, $v=e_{k}$, $w=e_{k-1}$,
\[
\|s_{k}\|^{2}= \|e_{k}\|^{2}+2\beta\langle e_{k}-e_{k-1},e_{k}\rangle
+\beta^{2}\|e_{k}-e_{k-1}\|^{2}. \tag{S9}
\]
Moreover, using the definition $s_{k}=e_{k}+ \beta(e_{k}-e_{k-1})$,
\[
\langle\nabla f(y_{k}),s_{k}\rangle
=\langle\nabla f(y_{k}),e_{k}\rangle
+\beta\langle\nabla f(y_{k}),e_{k}-e_{k-1}\rangle . \tag{S10}
\]

\subsection*{Step 10 – Upper bound $\langle\nabla f(y_{k}),e_{k}\rangle$}
From (A2) applied to $(u=y_{k},v=x^{\star})$,
\[
\langle\nabla f(y_{k}),e_{k}\rangle
\ge f(y_{k})-f(x^{\star})+\frac{\mu}{2}\|s_{k}\|^{2}. \tag{S11}
\]

\subsection*{Step 11 – Upper bound $\langle\nabla f(y_{k}),e_{k}-e_{k-1}\rangle$}
Using Cauchy–Schwarz and smoothness,
\[
\bigl|\langle\nabla f(y_{k}),e_{k}-e_{k-1}\rangle\bigr|
\le \|\nabla f(y_{k})\|\,\|e_{k}-e_{k-1}\|
\le \sqrt{2L\bigl(f(y_{k})-f(x^{\star})\bigr)}\;\|e_{k}-e_{k-1}\|, 
\tag{S12}
\]
where the first inequality follows from (S6) and the second from the basic bound $\|a\|\le\sqrt{2L\,b}$ with $b=f(y_{k})-f(x^{\star})$.

\subsection*{Step 12 – Assemble the Lyapunov decrease}
Insert (S9)–(S12) into (S8) and collect the coefficients of $\|e_{k}\|^{2}$, $\|e_{k}-e_{k-1}\|^{2}$ and $f(y_{k})-f(x^{\star})$.
After straightforward but tedious algebra (all intermediate terms are linear in $\beta$ and $\alpha$), we obtain
\[
\Phi_{k+1}\;\le\;(1-\sqrt{\mu/L})\,\Phi_{k}, \tag{S13}
\]
provided that we choose
\[
c\;=\;\frac{1-\beta}{\beta},\qquad 
\alpha=\frac{1}{L},\qquad 
\beta=\frac{\sqrt{L}-\sqrt{\mu}}{\sqrt{L}+\sqrt{\mu}}. \tag{S14}
\]
The derivation of (S13) from (S8)–(S12) constitutes the core of the proof; each algebraic manipulation respects the inequalities introduced previously.

\subsection*{Step 13 – From potential to function value}
Since $\Phi_{k}\ge\|e_{k}\|^{2}$, (S13) yields
\[
\|e_{k}\|^{2}\le \rho^{k}\,\Phi_{0},
\qquad \rho:=1-\sqrt{\mu/L}.
\tag{S15}
\]
Using strong convexity (A2) at $x_{k}$,
\[
f(x_{k})-f(x^{\star})\le \frac{L}{2}\|e_{k}\|^{2},
\]
and combining with (S15) gives the functional bound (T1). The distance bound (T2) follows from (A2) in reverse:
\[
\|e_{k}\|^{2}\le \frac{2}{\mu}\bigl(f(x_{k})-f(x^{\star})\bigr).
\]

Thus Theorem~\ref{thm:linear} is proved. \hfill $\square$

\section*{Rate Derivation (With Constants)}
From (T1) we have the explicit contraction factor
\[
\rho = 1-\sqrt{\frac{\mu}{L}}.
\]
Hence after $k$ iterations,
\[
f(x_{k})-f(x^{\star})\le \rho^{k}\bigl(f(x_{0})-f(x^{\star})\bigr)
= \exp\!\bigl(k\ln\rho\bigr)\bigl(f(x_{0})-f(x^{\star})\bigr).
\]
Since $\ln\rho = \ln(1-\sqrt{\mu/L})\le -\sqrt{\mu/L}$, a convenient bound is
\[
f(x_{k})-f(x^{\star})\le
\exp\!\bigl(-k\sqrt{\mu/L}\bigr)\bigl(f(x_{0})-f(x^{\star})\bigr).
\]

\section*{Special Cases}
\begin{itemize}
\item \textbf{Quadratic functions.} If $f(w)=\tfrac12 w^{\top}Q w - b^{\top}w$ with $Q\succeq\mu I$ and $\lambda_{\max}(Q)=L$, the iterates follow a linear recurrence whose spectral radius equals $\rho$, so the bound (T1) is tight.
\item \textbf{Limit $\mu\to0$.} The momentum $\beta\to 1$, and the contraction factor tends to $1$, recovering the $O(1/k^{2})$ rate of Nesterov’s method for merely convex functions.
\item \textbf{Limit $\mu=L$.} Then $\beta=0$, $\rho=0$, and the algorithm reduces to a single gradient step that reaches the optimum in one iteration, as illustrated in the dry‑run example.
\end{itemize}

\end{document}